\section{A Transformer instance}\label{sec:transformer-realization}

Nothing so far depends on attention. A standard Transformer block simply gives a familiar place to locate the represented cuts: source-resolved maps construct a receiver exposure, a coarser interface presents that exposure to a local continuation, and later blocks act on the represented state produced before them.

\subsection{Source-resolved additive routing}

Let $R$ be a finite receiver set and $S$ a finite source set. For source spaces $U_i$ and receiver exposure spaces $W_r$, let
\[
u_i:X\to U_i,
\qquad
C_{ri}:X\to\mathcal L(U_i,W_r).
\]

\begin{definition}[Additive routed exposure]\label{def:additive-routing-v13}
A source-resolved additive receiver exposure is
\[
\boxed{E_r(H)=\sum_{i\in S}C_{ri}(H)u_i(H).}
\]
The source-resolved predecessor state for that aggregation is
\[
P_r(H)=\bigl(C_{ri}(H)u_i(H)\bigr)_{i\in S}.
\]
\end{definition}

The composite exposure $E_r$ forgets a finer represented cut:
\[
X\xrightarrow{P_r}\mathcal M_r\xrightarrow{\sum}W_r.
\]
The kernel of the sum records which source-resolved states collide, while marked support records which sources can affect receiver $r$. A further relevance/payload factorization becomes architectural only when it is itself represented or marked.

Standard masked attention is one instance of this routed form. With receiver--source scores $s_{ri}=\langle q(x_r),k(y_i)\rangle/\tau$, admissible-source mask $S_r$, positive baselines $\kappa_{ri}$, and source-local values $v_i$, the familiar coefficients are
\[
\alpha_{ri}
=
\frac{\kappa_{ri}e^{s_{ri}}}{\sum_{j\in S_r}\kappa_{rj}e^{s_{rj}}}
\quad(i\in S_r),
\qquad
\alpha_{ri}=0\quad(i\notin S_r),
\]
and $E_r(H)=\sum_i\alpha_{ri}(H)v_i(H)$ \citep{vaswani2017attention}. Deriving this familiar formula from more primitive commitments requires assumptions about scalar evidence, positive multiplicative support, ratio-preserving row normalization, source-local values, and shared finite-rank scores; those assumptions are background to the represented-process comparison rather than part of its architecture identity criterion \citep{luce1959individual}.

\subsection{Multi-head attention as represented exposure}

For a nonempty finite head set $\mathcal H$, residual space $V$, and head maps
\[
Q_m,K_m:V\to\R^{d_m},
\qquad
V_m:V\to W_m,
\qquad
O_m:W_m\to V,
\]
let $N_1$ be the first receiver-local normalization and write the usual masked attention coefficients as $\alpha_{ij}^{(m)}(H)$.

\begin{proposition}[Multi-head source-resolved exposure]\label{prop:multihead-v13}
With optional receiver-local affine bias $b_{\mathrm{att}}\in V$,
\[
\boxed{
t_i(H)=b_{\mathrm{att}}+
\sum_{m\in\mathcal H}O_m
\left(\sum_{j\in S}\alpha_{ij}^{(m)}(H)V_mN_1(h_j)\right).}
\]
Hence $t_i(H)$ is the aggregate of the marked head--source contributions
\[
\bigl(\alpha_{ij}^{(m)}(H)O_mV_mN_1(h_j)\bigr)_{(m,j)}.
\]
\end{proposition}

\begin{proof}
The standard output projection decomposes into block maps $O_m$; linearity distributes each block over its source sum.
\end{proof}

At a source-resolved cut the head--source contributions are still represented separately. At the coarser receiver interface used in the contextual witness, only their aggregate $t_i(H)$ is retained alongside the receiver state $h_i$. These cuts answer different architectural questions.

\subsection{Receiver-local continuation and the FFN}

Let
\[
\FFN(x)=W_2\phi(W_1x+b_1)+b_2,
\qquad
a(x)=\phi(W_1x+b_1),
\]
with coordinatewise $\phi$. If $(e_e)$ is the standard basis of the hidden space and $w_e=W_2e_e$, then:

\begin{proposition}[Marked FFN hidden carrier]\label{prop:ffn-v13}
\[
\boxed{\FFN(x)=b_2+\sum_e a_e(x)w_e.}
\]
The specified continuation therefore factors through the marked hidden activation carrier $a(x)$.
\end{proposition}

The hidden activation is an actual represented state used by the FFN. Treating its coordinates as separate experts or mechanisms would require a stronger claim, and a hidden-basis mixing is passive only when its marks and incident maps are transported coherently.

\subsection{A standard PreNorm block}

Let $N_2:V\to V$ be the second receiver-local normalization and define
\[
u_i=h_i+t_i(H),
\qquad
h_i^+=u_i+\FFN(N_2(u_i)).
\]

\begin{theorem}[Receiver-process decomposition of a PreNorm block]\label{thm:prenorm-v13}
For receiver $i$, define
\[
Q_i^{\mathrm{att}}(H):=(h_i,t_i(H))
\]
corestricted to its realized image, and
\[
G_i^{\mathrm{block}}(h,t):=\nu+\FFN(N_2(\nu)),
\qquad
\nu=h+t.
\]
Then
\[
\boxed{B_i^{\mathrm{block}}=G_i^{\mathrm{block}}Q_i^{\mathrm{att}}.}
\]
The interface side admits the source-resolved head--source process of \Cref{prop:multihead-v13}; the continuation side admits the FFN hidden-carrier process of \Cref{prop:ffn-v13}.
\end{theorem}

\begin{proof}
Evaluating $G_i^{\mathrm{block}}$ at $(h_i,t_i(H))$ gives $u_i+\FFN(N_2(u_i))=h_i^+$.
\end{proof}

This is the concrete receiver process used earlier: the block presents $(h_i,t_i(H))$ and continues locally through the residual and FFN update. A finer cut can expose $U=(u_i)_i$ as another represented state. The overall block function does not choose which intermediate cut should count as the architecture under comparison.

\subsection{Depth as repeated process composition}

For blocks $B_0,\ldots,B_{L-1}$, let
\[
A_\ell:=B_{\ell-1}\circ\cdots\circ B_0,
\qquad A_0=\Id.
\]
A layer-$\ell$ receiver interface $Q_{\ell,j}$ has earlier-input distinction shadow $\mathsf D(Q_{\ell,j}A_\ell)$, family envelope $J_{A_\ell,\mathcal Q_{\ell,j}}$, and effective class set $\mathcal C_{A_\ell}(\mathcal Q_{\ell,j})$. Recutting gives these objects by repeated kernel pullback, while marked dependency neighborhoods compose by predecessor-neighborhood union. PreNorm, PostNorm, parallel branches, recurrence, adaptive depth, and carrier changes can therefore differ in the schedule of represented states on which later receivers are built even when a coarser formula-level description hides those differences.
